% introduction.tex
Functional MRI has become the workhorse of cognitive and systems
neuroscience over the past three decades. By detecting small changes
in blood oxygenation that accompany neural activity — the BOLD
(blood-oxygen-level-dependent) signal first described by Ogawa
et~al.~\cite{Ogawa1990} — fMRI lets us map which brain regions are
doing what, without any surgical intervention. What makes it
particularly useful is that the same scanner, with the same subject,
can probe entirely different cognitive systems simply by changing the
task.

However, raw fMRI data are messy. Every subject moves during the
scan. Slices of the brain are acquired at slightly different moments
in time. The functional and structural images are not naturally
aligned. Before any neuroscience can be done, the data need to go
through a preprocessing pipeline that corrects for these
acquisition-related artefacts and transforms all images into a common
reference space. Only after preprocessing can we ask statistical
questions about which regions activate in response to a given
stimulus~\cite{Poldrack2011}.

The standard tool for fMRI analysis in the neuroimaging community is
SPM (Statistical Parametric Mapping), developed at University College
London. SPM frames the whole problem as a general linear model
(GLM): each time series is regressed against a design matrix built
from the experimental conditions, convolved with a model of the
haemodynamic response~\cite{Friston1994}. The resulting parameter
estimates yield contrast images and t-statistics that localise
task-related activations. Beyond task-based analyses, the same
preprocessed data can be used to compute seed-based functional
connectivity — correlating the time series of a chosen reference
region against the rest of the brain~\cite{Biswal1995}.

This project implements the full pipeline — from raw NIfTI download
through preprocessing, GLM, and connectivity — as a set of seven
self-contained, sequentially runnable MATLAB scripts. Two datasets
from the OpenNeuro repository were chosen deliberately to cover
different functional systems: ds000114 targets the motor cortex
through a finger/foot/lips tapping task~\cite{Gorgolewski2013}, and
ds000105 targets the ventral visual stream through a
multi-category object-viewing experiment~\cite{Haxby2001}. Running
both datasets through the same pipeline allows us to validate
that the code produces anatomically sensible results for distinct
functional domains and to compare data quality metrics across the
two acquisitions.

The rest of this report is organised as follows. Section~II describes
the two datasets and their experimental designs. Section~III
(Methods) walks through every processing step with the parameters
used. Section~IV (Results) presents quality-control figures,
activation maps, and connectivity maps for both datasets together
with a cross-dataset comparison. Section~V discusses what the results
mean, and Section~VI concludes.
